\subsection{Quantum Approximate Optimization Algorithm}\label{sec:quantum-approximate-optimization-algorithm-qaoa}

\subsubsection{Why QAOA?}\label{sec:why-qaoa}

In the Variational Quantum Algorithms (VQAs) framework, the VQE and QAOA algorithms are \textbf{hybrid quantum-classical algorithms}. It means that the quantum computer executes a parameterized quantum circuit (ansatz), the classical computer updates its parameters, and the process repeats until convergence. \hl{QAOA is a Variational Algorithm for combinatorial optimization problems on NISQ devices and is a \textbf{special case} of VQE.}

\highspace
\begin{flushleft}
    \textcolor{Red2}{\faIcon{exclamation-triangle} \textbf{The problem}}
\end{flushleft}
Imagine we work in industry. Every week we receive a different optimization problem:
\begin{itemize}
    \item Routing trucks (e.g., UPS, FedEx)
    \item Scheduling workers (e.g., Amazon, Walmart)
    \item Assigning frequencies (e.g., AT\&T, Verizon)
    \item Portfolio optimization (e.g., BlackRock, Vanguard)
    \item Max-Cut (e.g., Facebook, Twitter)
    \item Maximum Independent Set (e.g., Google, Microsoft)
    \item SAT (e.g., IBM, Intel)
    \item Graph partitioning (e.g., Netflix, Spotify)
\end{itemize}
We could design a different quantum algorithm for each one of these problems, but it would be very inefficient. Instead, we want to design a \textbf{general-purpose quantum algorithm} that can solve all these problems. In other words, \textbf{one optimization framework that works for many problems}. This is exactly the goal of QAOA.

\highspace
\begin{flushleft}
    \textcolor{Green3}{\faIcon[regular]{lightbulb} \textbf{The idea}}
\end{flushleft}
Instead, industry wants something different. They want an algorithm that is:
\begin{itemize}
    \begin{multicols}{2}
        \item Generic
        \item Reusable
        \item Easy to adapt
        \item Requires minimal redesign
    \end{multicols}
\end{itemize}
QAOA takes an optimization problem, represents it as a Hamiltonian, and then run the same algorithm to find the solution. This is the fundamental idea of QAOA: the algorithm \textbf{does not change}; only the \textbf{Hamiltonian changes}.

\highspace
\textcolor{Green3}{\faIcon{check-circle}} This method has a great advantage: \hl{every combinatorial optimization problem can be represented in the same mathematical form.}

\highspace
\textcolor{Red2}{\faIcon{times-circle}} However, nothing comes for free. To make every optimization problem look identical, we must \textbf{remove all problem-specific information}. So, we \hl{lose the knowledge of the specific problem structure, which classical algorithms often exploit.} For example, classical Max-Cut algorithms know graphs, edges and connectivity; after converting everything into a Hamiltonian, QAOA only ``sees'' a Hamiltonian, it no longer knows what the vertices represent, why two variables interact, what the original optimization problem looked like. Everything becomes an \hl{energy minimization problem.}

\highspace
In summary, the \definition{Quantum Approximate Optimization Algorithm\break (QAOA)} is a \textbf{hybrid Variational Quantum Algorithm (VQA)} designed to \textbf{solve combinatorial optimization problems on NISQ quantum computers}. It is a special case of VQE. Rather than designing a different quantum algorithm for every optimization problem, QAOA converts the problem into a Cost Hamiltonian $H_C$ and applies the same optimization framework. Its main advantage is that many optimization problems can be expressed in the same Hamiltonian form, making the algorithm general and reusable. The drawback is that converting every problem into a generic Hamiltonian loses the original problem structure, which efficient classical algorithms often exploit.